\section{Methodology}
\label{sec:method}

In this section, we present the mathematical formulation of Rademacher-Bounded Fourier Trajectory Merging (RB-FTM). We first formalize weight-space task vector merging, then introduce the Fourier trajectory representation, prove its learning-theoretic complexity guarantees, and detail our optimization objective and its stability properties.

\subsection{Weight-Space Task Vector Merging}
Let $W_0^{(l)} \in \mathbb{R}^{D_l}$ denote the pre-trained parameter weights of a base neural network at layer $l \in \{0, \dots, L-1\}$, where $L$ is the total number of layers. Suppose we have $K$ distinct task-specific expert models $\{W_k\}_{k=0}^{K-1}$ that have been fine-tuned from the same shared pre-trained initialization $W_0$. The task-specific weight vectors (or task vectors) at layer $l$ are defined as:
\begin{equation}
    V_k^{(l)} = W_k^{(l)} - W_0^{(l)}
\end{equation}
For any given layer-wise ensembling coefficient matrix $\Lambda \in [0, 1]^{K \times L}$, where $\alpha_k(l) = \Lambda_{k,l}$ represents the ensembling coefficient for task $k$ at layer $l$, the merged model weights $W_{\text{merged}}^{(l)}(\Lambda)$ are assembled via:
\begin{equation}
    W_{\text{merged}}^{(l)}(\Lambda) = W_0^{(l)} + \sum_{k=0}^{K-1} \alpha_k(l) V_k^{(l)}
\end{equation}
Input representations propagate sequentially through the merged layers, yielding stable multi-task predictions when the layer-wise coefficients $\alpha_k(l)$ are smoothly and robustly configured.

\subsection{Fourier Trajectory Representation}
Rather than treating each element of $\Lambda$ as an independent parameter, which leads to an overparameterized search space of dimension $K \times L$, we parameterize the ensembling trajectory across depth as a continuous curve. 

We define the normalized network depth coordinates as $z = l / (L-1) \in [0, 1]$ for layer $l \in \{0, \dots, L-1\}$. 

We parameterize the ensembling coefficient $\alpha_k(l)$ for task $k$ at layer $l$ using a clipped continuous Fourier trajectory $T_k(z)$ of spectral cutoff frequency $F \ge 1$ ($F \ll L$):
\begin{equation}
    \alpha_k(l; \theta_k) = \Pi_{[0,1]} \left( T_k(z) \right)
\end{equation}
where $T_k(z)$ is the continuous Fourier series:
\begin{equation}
    T_k(z) = a_{k,0} + \sum_{f=1}^F \left( a_{k,f} \cos(2\pi f z) + b_{k,f} \sin(2\pi f z) \right)
\end{equation}
where:
\begin{itemize}
    \item $\Pi_{[0,1]}(x) = \max(0, \min(x, 1))$ is the projection (clipping) operator onto the interval $[0, 1]$.
    \item $\theta_k = [a_{k,0}, a_{k,1}, b_{k,1}, \dots, a_{k,F}, b_{k,F}]^T \in \mathbb{R}^{2F+1}$ represents the learned Fourier coefficients.
    \item $\Theta = \{\theta_k\}_{k=0}^{K-1} \in \mathbb{R}^{K \times (2F+1)}$ denotes the global continuous trajectory parameter matrix.
\end{itemize}
By initializing $a_{k,0} = 1/K$ and all harmonic coefficients $a_{k,f}, b_{k,f} = 0$, the optimization starts exactly at the robust static uniform merging baseline.

\subsection{Learning-Theoretic Complexity Guarantees}
To analyze the generalization capacity of our ensembling class, we formulate the hypothesis class of Fourier ensembling trajectories. Let $\mathcal{H}_F$ denote the class of Fourier ensembling trajectories of spectral cutoff $F$ with $\ell_1$-bounded coefficients, defined as $\mathcal{H}_F = \{ z \mapsto a_0 + \sum_{f=1}^F ( a_f \cos(2\pi f z) + b_f \sin(2\pi f z) ) \mid \|\theta\|_1 \le C_0 \}$, where $\|\theta\|_1 = |a_0| + \sum_{f=1}^F (|a_f| + |b_f|)$. 

The following theorem derives a mathematically rigorous bound on the empirical Rademacher complexity of $\mathcal{H}_F$ over network depth, proving that the ensembling capacity is strictly bounded by the low-frequency spectral cutoff $F$.

\begin{theorem}[Empirical Rademacher Complexity of Fourier Trajectories]
\label{thm:rademacher}
Let $\mathcal{H}_F$ be the class of Fourier ensembling trajectories of spectral cutoff $F$ with $\ell_1$-bounded coefficients $\|\theta\|_1 \le C_0$. The empirical Rademacher complexity of $\mathcal{H}_F$ over the network depth coordinate set $Z = \{z_1, \dots, z_L\}$ of size $L$ satisfies:
\begin{equation}
    \widehat{\mathcal{R}}_L(\mathcal{H}_F) \le C_0 \sqrt{\frac{2 \ln(4F+2)}{L}}
\end{equation}
\end{theorem}

\begin{proof}[Proof Sketch]
Since any Fourier trajectory $h(z)$ can be expressed as a linear combination of $2F+1$ bounded trigonometric basis functions scaled by $C_0$, we can express $\mathcal{H}_F$ as a subset of the convex hull of these basis functions. Applying H{\"o}lder's inequality and maximizing over the unique basis coordinate dimensions (accounting for positive and negative directions, yielding $4F+2$ elements), we apply Massart's Finite Lemma to obtain the precise logarithmic bound in terms of the spectral cutoff $F$ and network depth $L$. See Appendix~\ref{proof:thm_rademacher} for the complete, step-by-step mathematical proof.
\end{proof}

Crucially, because the clipping operator $\Pi_{[0,1]}(x)$ is $1$-Lipschitz, Talagrand's Contraction Lemma guarantees that the empirical Rademacher complexity of the clipped hypothesis class is upper bounded by that of the unclipped class, preserving the validity of the bound.

\subsection{Conceptual Distinction: Trajectory-Space vs. Data-Space Generalization}
We must highlight an important conceptual distinction regarding the learning-theoretic bound in Theorem~\ref{thm:rademacher}. In standard statistical learning theory, Rademacher complexity is evaluated over independent, identically distributed (i.i.d.) samples drawn from a data distribution to bound prediction generalization over unseen samples. Here, the "samples" are the network layer depth coordinates $Z = \{z_1, \dots, z_L\}$, which are fixed, deterministic, and highly ordered. 

Therefore, Theorem~\ref{thm:rademacher} does not bound the model's predictive generalization error on unseen data samples. Rather, it mathematically bounds the \emph{structural complexity of the trajectory coefficient function class itself}. By showing that $\widehat{\mathcal{R}}_L(\mathcal{H}_F)$ scales with the spectral cutoff frequency $F$ and decays with network depth $L$, we prove that the trajectory's capacity to fluctuate across depth is strictly bounded. This acts as a powerful structural regularizer over the weight assembly space, preventing the ensembling weights from overfitting to local representation noise on small calibration sets.

\subsection{Downstream Prediction Generalization Bridge}
We mathematically bridge our trajectory-space complexity bounds to the actual downstream prediction generalization error of the merged model on unseen data samples in Appendix~\ref{sec:appendix_generalization_bridge}. We first apply the vector-valued Ledoux-Talagrand Contraction Lemma to obtain a parameter-contraction bound, and discuss how the propagated Lipschitz constant $L_{\text{prop}}$ scales across the network depth $L$. 

We must explicitly address a critical theoretical limitation of standard composition-based generalization bounds in deep networks. Because ensembling coefficients enter multiplicatively across layers, the propagated Lipschitz constant $L_{\text{prop}}$ can scale exponentially with depth $L$ (i.e., $\mathcal{O}(\prod_{l=1}^L \Lambda_l)$ where $\Lambda_l$ is the Lipschitz constant of the $l$-th network block) unless the individual blocks are strictly contractive ($\Lambda_l \le 1$). In standard deep network architectures (such as CNNs and Transformers), blocks are typically not strictly contractive, as residual connections and self-attention mechanism weights are scaled to preserve representational norms. This exponential scaling can render standard composed bounds vacuous for very deep networks. In practice, however, architectural features such as Layer Normalization, Batch Normalization, and weight spectral normalization act as strong empirical regularizers that bound representation scales on a compact manifold, effectively preventing exponential representation explosion across depth. We discuss this limitation and the stabilizing role of normalization layers in Appendix~\ref{sec:appendix_generalization_bridge}.

To resolve the dimensional mismatch (the lack of a $1/\sqrt{N}$ sample size decay rate in standard contraction-based bounds), we provide a complete, formal derivation of a downstream prediction generalization bound via covering numbers of the trajectory-parameterized network class in Appendix~\ref{sec:appendix_generalization_bridge}. This formal covering-number derivation establishes an explicit $\widetilde{\mathcal{O}}(1/\sqrt{N})$ decay rate over data samples:
\begin{equation}
    \widehat{\mathcal{R}}_N(\mathcal{F}_{\mathcal{H}_F}) \le \mathcal{O}\left( \sqrt{\frac{d \ln(K C_0 L_{\Theta} N)}{N}} \right)
\end{equation}
where $d$ is the trajectory parameter dimension ($d = K(2F+1)$ for Fourier, and $d = K(F+1)$ for DCT), $C_0$ bounds the parameter norms, and $L_{\Theta}$ is the network's Lipschitz constant with respect to ensembling parameters. This theoretical bridge provides a rigorous guarantee: constraining the trajectory-space spectral complexity over depth coordinates strictly bounds the downstream prediction generalization error on unseen data samples. Crucially, while this bound is completely independent of the parameter count of the underlying neural network (mathematically proving that spectral trajectory optimization avoids the curse of model parameter overparameterization), the trajectory parameter dimension $d$ scales linearly with $K$. Consequently, we transparently disclose that the actual model's predictive generalization capability scales as $\mathcal{O}(\sqrt{K/N})$, meaning that the ensembling capacity of the actual merged network does depend on the number of task experts $K$.

\subsection{Multi-Task Joint Trajectory Complexity Scaling}
We also derive joint multi-task ensembling trajectory complexity bounds. We first present a stylized joint formulation where the complexity of the joint multi-task trajectory class $\mathcal{H}_F^{\text{joint}}$ is strictly independent of the task count $K$ inside the logarithm.
\begin{theorem}[Joint Multi-Task Trajectory Complexity]
\label{thm:rademacher_joint}
Let $\mathcal{H}_F^{\text{joint}}$ be the joint multi-task trajectory class for $K$ tasks ensembled simultaneously. The empirical Rademacher complexity of this joint class over network depth coordinates of size $L$ is bounded independently of the task count $K$:
\begin{equation}
    \widehat{\mathcal{R}}_L(\mathcal{H}_F^{\text{joint}}) \le C_{\text{joint}} \sqrt{\frac{2 \ln(4F+2)}{L}}
\end{equation}
For the joint Discrete Cosine (DCT) multi-task trajectory class $\mathcal{H}_F^{\text{joint, DCT}}$, the joint complexity satisfies:
\begin{equation}
    \widehat{\mathcal{R}}_L(\mathcal{H}_F^{\text{joint, DCT}}) \le C_{\text{joint}} \sqrt{\frac{2 \ln(2F+2)}{L}}
\end{equation}
\end{theorem}
We present the proof sketch and complete joint derivation in Appendix~\ref{proof:thm_rademacher_joint}.

\begin{remark}[Shared vs. Independent Rademacher Variables and the Task-Scaling Discrepancy]
In Theorem~\ref{thm:rademacher_joint}, our joint trajectory is evaluated as a stylized, scalar-valued sum $g(z) = \sum_k \alpha_k(z)$, employing a single set of Rademacher variables $\sigma_l$ across depth coordinates. While this abstraction mathematically eliminates $K$ from the trajectory-space bound, we emphasize that the actual merged neural network does not depend on a simple scalar sum of trajectories. Because task-specific expert vectors $\{V_k^{(l)}\}$ are geometrically distinct and operate in independent weight-space directions, downstream predictions depend on each ensembling trajectory $\{\alpha_k(l)\}$ independently. Consequently, as established in the downstream generalization bound (Eq. 16), the actual predictive capacity scales as $\mathcal{O}(\sqrt{K/N})$.

For completeness, if we define independent Rademacher variables $\sigma_{l,k}$ for each task $k$ and layer $l$ under a standard vector-valued multi-task complexity framework, we obtain a bound that explicitly depends on $K$:
\begin{align}
    \widehat{\mathcal{R}}_L(\boldsymbol{\mathcal{H}}_F^{\text{multi}}) &\le C_{\text{joint}} \sqrt{\frac{2 \ln(2K(2F+1))}{L}} \nonumber \\
    &= C_{\text{joint}} \sqrt{\frac{2 \ln(4KF+2K)}{L}}
\end{align}
We provide a comprehensive discussion and formal derivation of this independent-variable multi-task bound in Appendix~\ref{proof:thm_rademacher_joint_independent}. Both formulations show that the scaling with respect to task count $K$ is either completely independent (under the scalar projection) or strictly logarithmic ($\mathcal{O}(\sqrt{\ln(KF)/L})$ under the vector-valued class) in trajectory space. This ensures that the trajectory structural complexity scales exceptionally well as the number of tasks $K$ grows.
\end{remark}

\subsection{Clipping Projections and Spectral Leakage Analysis}
To restrict the ensembling coefficients to the valid range $[0, 1]$, we apply the $1$-Lipschitz clipping projection operator, which preserves our learning-theoretic complexity bounds. While hard clipping represents a non-linear thresholding operation that introduces high-frequency harmonic overtones (spectral leakage) in the Fourier decomposition, we show in Appendix~\ref{sec:appendix_clipping_leakage} that the leaked energy decays algebraically as $\mathcal{O}(1/n^2)$ for the $n$-th harmonic overtone. Furthermore, because representation propagation acts as a recurrent update (effectively a discrete-time integrator), it behaves as a natural low-pass filter that suppresses these high-frequency components, keeping the activations smooth across depth.

\subsection{Optimization Objective and Spectral Lasso}
To physically enforce the learning-theoretic complexity bound derived in Theorem~\ref{thm:rademacher}, we directly couple our theoretical constraint to a practical regularized loss function. We formulate the \emph{Spectral-Regularized Few-Shot Calibration} objective:
\begin{equation}
    \min_{\Theta} \mathcal{L}(\Theta) = \mathcal{L}_{\text{CE}}(\mathcal{D}_{\text{cal}}; \Theta) + \gamma \sum_{k=0}^{K-1} \|\theta_{k,\text{harm}}\|_1
\end{equation}
where $\mathcal{L}_{\text{CE}}$ is the Cross-Entropy loss evaluated on a small calibration dataset $\mathcal{D}_{\text{cal}}$, and $\gamma \ge 0$ is the spectral regularization coefficient. 

Crucially, the $L_1$ penalty (Spectral Lasso) is applied strictly to the harmonic coefficients, denoted by $\theta_{k,\text{harm}} = [a_{k,1}, b_{k,1}, \dots, a_{k,F}, b_{k,F}]^T$ (for Fourier trajectories) and $\theta^{\text{DCT}}_{k,\text{harm}} = [a_{k,1}, \dots, a_{k,F}]^T$ (for DCT trajectories), while excluding the base uniform coefficient $a_{k,0}$. Because $a_{k,0}$ determines the baseline ensembling weight (initialized to the uniform value $1/K$), penalizing it towards zero would shrink the overall scale of the ensembling weights $\sum_k \alpha_k(l)$ across layers. This would artificially attenuate the representation propagation scale, leading to under-scaled activations at later layers and degrading final classification performance. 

By applying the Spectral Lasso strictly to the harmonic coefficients, the objective $\|\theta_{k,\text{harm}}\|_1 = \sum_{f=1}^F (|a_{k,f}| + |b_{k,f}|) \le C_0$ directly restricts the trajectory's capacity to fluctuate across network depth (forcing the trajectory back towards a flat, stable, and highly robust uniform baseline) while preserving the optimal, baseline activation scale. This limits the active ensembling capacity, preventing transductive overfitting on the few-shot calibration data while maintaining perfect representational hygiene.

\begin{remark}[Spectral Lasso ($L_1$) vs. Spectral Ridge ($L_2$)]
While we employ an $L_1$ Spectral Lasso penalty strictly on the harmonic coefficients ($\theta_{k,\text{harm}}$), one could alternatively formulate a quadratic $L_2$ Spectral Ridge penalty: $\gamma \sum_{k=0}^{K-1} \|\theta_{k,\text{harm}}\|_2^2$. The choice of $L_1$ regularization is theoretically and practically superior in our setting. First, the $L_1$ ball constraint matches the formulation of our Rademacher complexity bounds (Theorem~\ref{thm:rademacher}), which naturally employ an $\ell_1$-norm bound on $\theta$ to bound trajectory fluctuations via Massart's Finite Lemma. Second, $L_1$ regularization promotes sparsity, effectively acting as an automated spectral selection mechanism that prunes high-frequency harmonic overtones entirely. In contrast, an $L_2$ penalty would shrink all frequencies uniformly without setting any to zero, failing to discard unnecessary high-frequency capacity and resulting in a denser, less interpretable spectral trajectory.
\end{remark}

\begin{remark}[Theory-Practice Gap of the $L_1$ Penalty]
We acknowledge a standard machine learning theory-practice gap: while the Rademacher complexity bounds in Theorem~\ref{thm:rademacher} and Theorem~\ref{thm:rademacher_dct} are derived assuming a hard constraint on the parameter norm ($\|\theta\|_1 \le C_0$), our practical optimization objective employs a soft Lagrangian regularization penalty ($\gamma \sum_k \|\theta_{k,\text{harm}}\|_1$). By Lagrangian duality, for any choice of regularizer $\gamma \ge 0$, there exists an equivalent constraint radius $C_0 \ge 0$. However, the exact radius $C_0$ is data-dependent and is not explicitly bounded or quantified during optimization. This soft-regularization-vs-hard-constraint formulation is a standard and widely accepted practice in statistical learning optimization.
\end{remark}

\begin{remark}[Practical Selection and Heuristics for the Regularization Parameter $\gamma$]
In realistic, extremely resource-constrained few-shot calibration regimes (e.g., 10-shot adaptation), a large validation set is unavailable, making standard hyperparameter sweeps prone to severe transductive overfitting. To resolve this practical bottleneck, we propose two data-driven heuristics to select the spectral Lasso penalty $\gamma$ automatically:
\begin{enumerate}
    \item \textbf{Few-Shot Split-Validation}: Split the 10-shot calibration dataset into a 7-shot training split and a 3-shot validation split. Because our low-dimensional spectral optimization completes in under 15 seconds (requiring only 30 optimization steps of the Adam optimizer on 6 parameters), we run a rapid 3-step grid search over $\gamma \in \{10^{-3}, 10^{-2}, 10^{-1}\}$ inside 5 seconds, selecting the value that minimizes validation cross-entropy.
    \item \textbf{Spectral Covariance Heuristic}: Initialize $\gamma$ dynamically based on the pre-merging task vector scale relative to the base model. Let $\gamma_{\text{heuristic}} = \kappa \cdot \frac{1}{L} \sum_{l=0}^{L-1} \left( \frac{1}{K} \sum_{k=0}^{K-1} \|W_k^{(l)} - W_0^{(l)}\|_F^2 \right)$ where $\| \cdot \|_F$ is the Frobenius norm and $\kappa = 10^{-4}$ is a scaling constant. This dynamically scales up regularization when merging highly divergent experts to counteract representation shearing.
\end{enumerate}
\end{remark}

\subsection{Boundary Properties, Discrete Cosine Transform, and Tighter Complexity Bounds}
By adopting a periodic integer-frequency Fourier basis, our formulation mathematically establishes a periodic boundary identity at $z = 0$ (the first layer) and $z = 1$ (the last layer), forcing identical ensembling weights:
\begin{equation}
    \alpha_k(0) = \alpha_k(L-1) = a_{k,0} + \sum_{f=1}^F a_{k,f}
\end{equation}
Forcing identical weights at the first layer (low-level feature extraction) and the last layer (high-level classification projection) represents an artificial architectural constraint. To resolve this boundary identity while preserving the benefits of smooth spectral trajectories, we introduce the \textbf{Discrete Cosine Trajectory Merging (RB-DCTM)} variant, which utilizes a half-period cosine basis:
\begin{equation}
    \alpha^{\text{DCT}}_k(l; \theta_k) = \Pi_{[0,1]} \left( a_{k,0} + \sum_{f=1}^F a_{k,f} \cos\left(\pi f z\right) \right)
\end{equation}
Because the DCT basis uses half-frequencies ($\pi f z$ instead of $2\pi f z$), it eliminates the periodic identity, allowing the ensembling weights at the first layer ($z=0$) and last layer ($z=1$) to be optimized independently, while still maintaining the natural smoothness and boundary stability of harmonic functions. This also resolves the boundary runaway of polynomial trajectories (such as RBPM, $d=2$) caused by rigid shape constraints on $z \in [0, 1]$.

\begin{remark}[Homogeneous Neumann Boundary Conditions and Flat Derivatives]
While the half-period cosine basis $\cos(\pi f z)$ grants boundary value independence ($\alpha_k(0) \neq \alpha_k(1)$), it introduces an implicit homogeneous Neumann boundary condition on the derivatives of the trajectory, forcing $h'(0) = h'(1) = 0$ for any choice of coefficients. Architecturally, this flat-derivative constraint is highly beneficial: it guarantees that the ensembling weights remain stable and flat in the critical early layers (feature extraction) and late layers (classification projection), preventing rapid layer-wise fluctuations where the representational topology is most sensitive. Physically, this implicit flat-derivative constraint prevents high-frequency gradient updates from propagating into the boundary layers during few-shot optimization, creating a "boundary buffer" that protects the delicate initial representation extraction and final classification boundaries from destructive interference. While alternative bases like Chebyshev polynomials or splines would allow non-zero boundary derivatives, they also restore the risk of boundary runaway and high-frequency oscillations. The Neumann constraint of the DCT basis thus acts as an elegant, physics-inspired regularizer that stabilizes the boundaries.
\end{remark}

\begin{remark}[The Representational Flexibility Trade-Off of Neumann Boundary Conditions]
Although the homogeneous Neumann boundary condition $h'(0) = h'(1) = 0$ acts as a beneficial stabilizing buffer, we explicitly disclose that it represents a double-edged architectural constraint that restricts representational flexibility. In highly heterogeneous transfer learning scenarios where expert models are fine-tuned on extremely divergent input domains (e.g., merging a satellite image expert and a medical image expert), the first feature-extraction layers may need to undergo rapid, high-frequency transformations to resolve conflicting domain-specific features. Forcing a flat derivative at $z=0$ restricts the network's ability to transition ensembling weights rapidly at the input layer, potentially leading to a sub-optimal representational compromise. In such highly disparate settings, practitioners can relax the Neumann boundary condition by using a shifted Cosine basis, or by combining low-frequency sine and cosine harmonics (mixed Fourier-Cosine representation) to allow non-zero boundary derivatives while preserving low-frequency smoothness.
\end{remark}

Remarkably, RB-DCTM also yields a strictly tighter learning-theoretic complexity bound than its Fourier counterpart. Let the class of DCT ensembling trajectories $\mathcal{H}_F^{\text{DCT}}$ with $\ell_1$-bounded coefficients be defined as $\mathcal{H}_F^{\text{DCT}} = \{ z \mapsto a_0 + \sum_{f=1}^F a_f \cos(\pi f z) \mid \|\theta\|_1 \le C_0 \}$, where $\|\theta\|_1 = |a_0| + \sum_{f=1}^F |a_f|$. We now derive its empirical Rademacher complexity:

\begin{theorem}[Empirical Rademacher Complexity of DCT Trajectories]
\label{thm:rademacher_dct}
The empirical Rademacher complexity of the Discrete Cosine trajectory class $\mathcal{H}_F^{\text{DCT}}$ over the network depth coordinate set $Z = \{z_1, \dots, z_L\}$ of size $L$ satisfies:
\begin{equation}
    \widehat{\mathcal{R}}_L(\mathcal{H}_F^{\text{DCT}}) \le C_0 \sqrt{\frac{2 \ln(2F+2)}{L}}
\end{equation}
\end{theorem}

\begin{proof}[Proof Sketch]
By utilizing a cosine-only basis, the dimension of the basis vector $\Psi^{\text{DCT}}(z)$ is reduced from $2F+1$ to $F+1$. Maximizing over these components in H{\"o}lder's inequality corresponds to a finite set of sub-Gaussian random variables of size exactly $2(F+1) = 2F+2$ (accounting for positive and negative directions). Applying Massart's Finite Lemma yields a strictly tighter complexity bound than the Fourier alternative. See Appendix~\ref{proof:thm_rademacher_dct} for the complete, step-by-step mathematical proof.
\end{proof}

This tighter bound demonstrates that by omitting the sine harmonics and adopting a cosine-only basis, RB-DCTM reduces the trajectory space's structural complexity by a factor of inside the logarithm, providing a superior theoretical guarantee of generalization across layers while simultaneously granting boundary independence.

In our Analytical Coordinate Sandbox, representation propagation begins with the initial state vector $h_0 = X$ at layer $l = 0$. Since layer $0$ represents the unpropagated raw feature inputs, the propagation update equations naturally start at layer $l = 1$ and continue up to $l = L-1$. Under this framework, the first layer's ensembling parameters (`alpha[0, :]`) are structurally mapped to the initial coordinates, meaning propagation naturally begins at layer $1$ in the simulation code. This aligns the mathematical coordinate boundaries of representation propagation with the simulation implementation.